% Transcription — fonds Alexandre Grothendieck, shelfmark 57, batch 15 (pages 281-300).
% Established from the facsimile archives/batches/57/batch-15.pdf, per the skill
% transcribe-grothendieck.
%
% Pass: Opus 5.5 (claude-opus-5-5), 2026-10-03 — first pass, unchecked against the pages by a human.
%
% Pages transcribed: 282-291, 293, 296, 297, 299. Absent: 281 (cover sheet in
% his hand, its words used as the first heading), 292 (blank), 294, 295, 298, 300
% (unrelated typescripts: a Bourbaki-type text, a third-party typescript on lim^(p),
% two pages of a "Tribu" report).
%   282-289: examples for Pic: devissage of Pic for X -> Y -> S; quotients
%     (A x B)/Gamma of abelian schemes (Igusa-type example); p_* Omega^1;
%     a homomorphism divisible by p on the special fibre only (NS sections).
%   290-291: structure of p-primary divisible groups, T_p and T^p.
%   293, 296: notes on deformations and flatness of Pic, Pic of a quotient.
%   297, 299: counterexamples for gluing (two planes glued along a punctured line).
\documentclass[11pt,a4paper]{article}
\input{../preamble/grothendieck}

\folder{57}
\batch{15}
\pages{281}{300}
\foldertitle{Picard : tapuscrit annoté (s.d.), lettres (s.d., 1962).}
\dating{1962-[vers 1968]}
\watermark{Édition de démonstration}

\begin{document}

\section{Exemples}
\note{titre de sa main, seul mot d'une feuille par ailleurs blanche (p.~281)}

\subsection{Dévissage de $\underline{\mathrm{Pic}}$ pour un composé $X \to Y \to S$}

\page{282}
\marginal{\uncertain{L'exemple} \uncertain{produits}}

\begin{tikzcd}[row sep=small]
X \arrow[d, "g"] \\
Y \arrow[d, "f"] \\
S
\end{tikzcd}

\[
0 \to \mathrm{Pic}(Y) \to \mathrm{Pic}(X) \to \mathrm{Pic}(X/Y)
\qquad \mathrm{Pic}(X/Y) = \Gamma(\underline{\mathrm{Pic}}_{X/Y}/Y)
\]
\note{sous $\mathrm{Pic}(X/Y)$, un second signe $=$ renvoie à une expression biffée, illisible, où se lit seulement $\underline{\mathrm{Pic}}$ indicé par $Y/S$}
surjectif si $X/Y$ a une section.

Pourvu que $g$ soit qu.-cpt séparé et $\mathcal{O}_Y \xrightarrow{\sim} g_*(\mathcal{O}_X)$. Si cette dernière condition \uncertain{est} $S$-universelle, \uncertain{alors} pour faire un dévissage de \ill{} \uncertain{quelconque} $S' \to S$, et pour \ill{} des \ill{} \ill{} « faisceaux » courants \ill{} suite exacte de \struck{\ill{}} « préfaisceaux » \ill{} \ill{} une suite exacte
\[
0 \to \underline{\mathrm{Pic}}_{Y/S} \to \underline{\mathrm{Pic}}_{X/S} \to \textstyle\prod_{Y/S} \underline{\mathrm{Pic}}_{X/Y} ,
\]
\uncertain{ainsi} \uncertain{obtenue}, \ill{} \ill{} avec d'ailleurs une \uncertain{suite} \add{\uncertain{et} \uncertain{exacte}} à droite si $X/Y$ a une section. Dans ce cas les \uncertain{foncteurs} \ill{} \uncertain{représentables}, \ill{} \ill{} \ill{}
\[
0 \to \underline{\mathrm{Pic}}_{Y/S} \to \underline{\mathrm{Pic}}_{X/S} \to \textstyle\prod_{Y/S} \underline{\mathrm{Pic}}_{X/Y}
\]
\uncertain{où} la dernière flèche a une \uncertain{section} : \ill{} \ill{} si $X/Y$ a une section. Dans ce cas, on a donc
\[
\underline{\mathrm{Pic}}_{X/S} \simeq \underline{\mathrm{Pic}}_{Y/S} \times \textstyle\prod_{Y/S} \underline{\mathrm{Pic}}_{X/Y}
\]
[Ceci s'applique en particulier, \uncertain{comme} \ill{}, quand $X = Y \times_S Z$, $Z$ muni d'une section sur $S$ :
\[
\underline{\mathrm{Pic}}_{X/S} \simeq \underline{\mathrm{Pic}}_{Y/S} \times \underline{\mathrm{Pic}}_{Z/S} \times \underline{\mathrm{Hom}}_{\text{pointés}}(Y, \underline{\mathrm{Pic}}_{Z/S}) \;]
\]
\note{le dernier $\underline{\mathrm{Pic}}_{Z/S}$ porte un petit exposant, peut-être $\tau$, qui n'est pas sûr}

\subsection{Quotients $(A \times_S B)/\Gamma$}

\page{283}
Soit $S$ un préschéma de base, $A$ un schéma abélien sur $S$, $B$ un schéma \ill{} sur $S$, $\Gamma$ un groupe fini, opérant sur $A$ par automorphismes et sur $B$ \ill{}, \ill{} $\Gamma$ opère \uncertain{diagonalement} sur $A \times B$ \ill{} \ill{} \ill{} \ill{}, et \ill{} \ill{}
$X = (A \times B)/\Gamma$ comme \ill{} \ill{} \ill{} \uncertain{abélien} sur $Y = B/\Gamma$. On a donc, si \struck{\ill{}} $X/Y$ \ill{} une section,
\[
\underline{\mathrm{Pic}}_{X/S} \simeq \underline{\mathrm{Pic}}_{Y/S} \times \textstyle\prod_{Y/S} \underline{\mathrm{Pic}}_{X/Y} ,
\]
\drawing{dans la marge de gauche, le diagramme du quotient, reproduit ci-dessous ; un trait oblique y relie aussi $B$ à $S$}

\begin{tikzcd}
X \arrow[d, "f"'] & A \times_S B \arrow[l, "\Gamma"'] \arrow[d, "\mathrm{pr}_1"] \\
Y \arrow[d, "g"'] & B \arrow[l, "\Gamma"'] \\
S &
\end{tikzcd}

\note{l'étiquette de la flèche verticale de droite se lit $\mathrm{pr}_1$, alors qu'il s'agit de la projection $A \times_S B \to B$ ; on la garde telle quelle}
si \uncertain{($B \to S$} \ill{}\uncertain{)}, \ill{} disons \ill{} \ill{}, et $A \to S$ projectif, on aura
\[
\underline{\mathrm{Pic}}_{X/S} \simeq \underline{\mathrm{Pic}}_{Y/S} \times \textstyle\prod_{Y/S} \underline{\mathrm{Pic}}_{X/Y} ,
\qquad
\textstyle\prod_{Y/S} \underline{\mathrm{Pic}}_{(A' \times B)/\Gamma} \simeq \underline{\mathrm{Hom}}_{S,\Gamma}(B, \underline{\mathrm{Pic}}_{A/S})
\]
\note{la seconde formule est écrite en deux étages, l'accolade du haut portant $\prod_{Y/S} \underline{\mathrm{Pic}}_{(A' \times B)/\Gamma} \simeq$ et celle du bas $\simeq \underline{\mathrm{Hom}}_{S,\Gamma}(B, \underline{\mathrm{Pic}}_{A/S})$ ; le $A'$ dans la première est sa lecture, peut-être un lapsus pour $A$}

Soit $A'$ le schéma abélien dual de $A$, \uncertain{sur} \uncertain{lequel} opère $\Gamma$ par « transport de structure », $A'^{\Gamma}$ \emph{est} \ill{} \emph{à fibres discrètes} \ill{} \ill{} \ill{} \ill{} \ill{} \ill{} \ill{} \ill{} \ill{} sur $S$ [\ill{} \ill{} \ill{} \ill{} les opérations de $\Gamma$ \ill{} \ill{} $x \mapsto -x$] \emph{\uncertain{sur} $B$ \uncertain{non} \ill{}}. \struck{\ill{} \ill{} \ill{} \ill{} \ill{}} \ill{} \ill{} \ill{} \ill{} \uncertain{écrit}
\[
\underline{\mathrm{Hom}}_{S,\Gamma}(B, A')^{(\tau)} \simeq A'^{\Gamma} ,
\]
d'où \ill{}

\marginal{écrit en travers de la marge de gauche : et $\underline{N}^{\Gamma}$ \ill{} schéma \ill{} \uncertain{fini}, on a une suite exacte $0 \to \underline{\mathrm{Hom}}_{S,\Gamma}(B, A') \to \underline{\mathrm{Hom}}_{S,\Gamma}(B, \underline{\mathrm{Pic}}_{A/S})$ $\to \underline{\mathrm{Hom}}_{S,\Gamma}(B, \underline{\mathrm{NS}}_{A/S}) \simeq \underline{N}^{\Gamma}$, d'où $\underline{\mathrm{Pic}}^{\tau}_{X/S} \simeq \underline{\mathrm{Pic}}^{\tau}_{Y/S} \times \underline{\mathrm{Hom}}_{S,\Gamma}(B, A')^{\tau}$}

\page{284}
\[
\underline{\mathrm{Pic}}^{\tau}_{X/S} \simeq \underline{\mathrm{Pic}}^{\tau}_{Y/S} \times A'^{\Gamma}
\]

Par exemple, prenons $\Gamma = \mathbb{Z}/2\mathbb{Z}$, \add{opérant sur $A$ par \uncertain{symétrie},} $B$ un autre schéma abélien, sur lequel $\Gamma$ opère par translation par une section qui est partout d'ordre 2 exactement (i.e.\ \ill{} \ill{} une isom.\ de $\Gamma_S \to B$) --- donc $Y$ est également un schéma abélien sur $S$ --- alors
\[
\underline{\mathrm{Pic}}^{\tau}_{X/S} \simeq Y' \times {}_2A'
\]
On notera que $\underline{\mathrm{Pic}}^{\tau}_{X/S}$ est \uncertain{plat} sur $S$. En caractéristiques résiduelles $\neq 2$, $\underline{\mathrm{Pic}}^{\tau}_{X/S}$ est \struck{\uncertain{lisse}} \emph{\uncertain{simple}} sur $S$. Mais il ne l'est pas en un pt $s \in S$ de car.\ résiduelle 2. C'est l'exemple de \emph{\uncertain{Igusa}} [N.B. On peut prendre pour $A$, $B$ des \emph{courbes elliptiques}, \ill{} \ill{} \emph{\uncertain{surfaces}}. \struck{Cette \ill{}} \struck{\ill{}, on peut le supposer} On notera \emph{que $H^1(X, \mathcal{O}_X)$ est de dimension} \uncertain{égale} :
\[
\dim B + \dim A
\]
alors que $\dim \underline{\mathrm{Pic}}_{X/K} = \dim B$. On retrouve que l'exemple \uncertain{ci-dessus} donne \uncertain{lieu} à des « sauts » pour le $H^1(X, \mathcal{O}_X)$ \ill{} \uncertain{égales} \uncertain{caractéristiques}. Dans ces exemples, $\underline{\mathrm{Pic}}^{\tau}_{X/S}$ \ill{}

\page{285}
\uncertain{également} \ill{} « \uncertain{Albanese} » \ill{}, \ill{} schéma \uncertain{abélien}.

[\add{N.B.} Je n'étais pas \uncertain{certain} de \ill{}, \ill{} l'exemples qu'on \ill{} la \ill{} que $B$ soit un schéma \ill{} sur lequel $\Gamma$ opère \uncertain{par} \uncertain{translations}. En effet, on a une suite exacte
\[
0 \to \underline{\mathrm{Pic}}_{A/S} \to \underline{\mathrm{Hom}}_S(B, \underline{\mathrm{Pic}}_{A/S}) \to \underline{\mathrm{Hom}}_{S\text{-gr}}(A, \underline{\mathrm{Pic}}^{\tau}_{B/S}) \to 0
\]
\note{sous $\underline{\mathrm{Pic}}_{A/S}$ deux mots, peut-être « \uncertain{applications} \uncertain{constantes} » ; sous le dernier terme, une accolade : « [\uncertain{correspondances} \uncertain{divisorielles} sur $A \times B$] »}
\marginal{\ill{} \uncertain{relations} \uncertain{usitées} sur $B$ …}
d'où \uncertain{une} suite exacte
\[
0 \to \underline{\mathrm{Pic}}_{A/S}^{\Gamma} \to \underline{\mathrm{Hom}}_{S,\Gamma}(B, \underline{\mathrm{Pic}}_{A/S}) \to \underline{\mathrm{Hom}}_{S\text{-gr}}(A, \underline{\mathrm{Pic}}^{\tau}_{B/S})^{\Gamma}
\]
(\emph{et $\Gamma$ opère trivialement sur $\underline{\mathrm{Pic}}^{\tau}_{B/S}$}). Or \struck{\ill{}} supposons que $A_{\Gamma} = 0$ \uncertain{ensemblistement}, \uncertain{alors}
\[
\underline{\mathrm{Hom}}_{S\text{-gr}}(A, \underline{\mathrm{Pic}}^{\tau}_{B/S})^{\Gamma} = 0
\]
[utilisant le fait que \ill{} \ill{} \ill{} \uncertain{vérifié} sur $S$], i.e.\ un $S$-schéma \ill{} et il vient
\[
\underline{\mathrm{Pic}}_{A/S}^{\Gamma} \simeq \underline{\mathrm{Hom}}_{S,\Gamma}(B, \underline{\mathrm{Pic}}_{A/S}) .
\]
Donc \uncertain{on} \uncertain{aura}
\[
(\ast)\qquad
\begin{cases}
\underline{\mathrm{Pic}}_{X/S} \simeq \underline{\mathrm{Pic}}_{Y/S} \times \underline{\mathrm{Pic}}_{A/S}^{\Gamma} \\
\underline{\mathrm{Pic}}^{\tau}_{X/S} \simeq \underline{\mathrm{Pic}}^{\tau}_{Y/S} \times A'^{\Gamma}
\end{cases}
\]
\note{le crochet fermant qui suit est surchargé de traits}
On \uncertain{peut} \uncertain{remarquer} que \ill{} \ill{} \ill{} \ill{}, \ill{} bien que $A$ soit un schéma abélien, \uncertain{que} $\mathrm{Alb}(A)$ existe, et $(\mathrm{Alb}(A))_{\Gamma} = 0$ ensemblistement ….]

\subsection{Calcul de $p_*\underline{\Omega}^1_{X/S}$}

\page{286}
\begin{tikzcd}
X \arrow[d, "p"'] & A \times_S B \arrow[l, "\Gamma"'] \arrow[dl, "q"] \\
S &
\end{tikzcd}

On aura
\[
p_*(\underline{\Omega}^1_{X/S}) = q_*(\underline{\Omega}^1_{A \times_S B})^{\Gamma} = \pi_{A*}(\underline{\Omega}^1_{A/S})^{\Gamma} \times \pi_{B*}(\underline{\Omega}^1_{B/S})^{\Gamma}
\]
Dans le cas plus haut où $A$ et $B$ sont des schémas abéliens, \struck{on trouve} $\Gamma$ opérant \uncertain{sur} $B$ par translations, on trouve
\[
\text{a)}\qquad \omega^1_{X/S} \simeq \omega^1_{B/S} \times {\omega^1_{A/S}}^{\Gamma}
\]
Si $S$ est le spectre d'un corps $k$, pour comparer \struck{\ill{}} $\operatorname{rang} H^1(X, \mathcal{O}_X)$ et \struck{\ill{}} $\operatorname{rang}_k H^0(X, \Omega^1_{X/k})$, on note \struck{doit} que \ill{} \ill{} \ill{} \uncertain{on} \uncertain{aura}
\[
\begin{aligned}
&\text{a}')\quad H^0(X, \Omega^1_X) \simeq H^0(B, \Omega^1_B) \times H^0(A, \Omega^1_A)^{\Gamma} \\
&\text{b}')\quad H^1(X, \mathcal{O}_X) \simeq H^1(Y, \mathcal{O}_Y) \times \underbrace{H^1(A, \mathcal{O}_A)^{\Gamma}}_{t_{A'}^{\Gamma}}
\end{aligned}
\]
\uncertain{donc} \uncertain{comparant} \ill{} \ill{} \ill{}, on aura
\[
\operatorname{rang} H^1(X, \mathcal{O}_X) - \operatorname{rang} H^0(X, \Omega^1_{X/k}) = \dim t_{A'}^{\Gamma} - \dim (t_A)'^{\Gamma}
\]
Nous voulons donner des exemples \uncertain{où} \uncertain{cette} \uncertain{différence} soit $>0$, ou $<0$, \struck{et} ou \ill{} [dans le cas d'une \uncertain{famille}] \struck{dans} \add{les} termes $\dim (t_{A'})^{\Gamma}$ et $\dim (t_A)'^{\Gamma} =$ \struck{$\dim$} $\dim (t_A)_{\Gamma}$ variant d'une valeur \uncertain{en} \uncertain{particulier} à l'autre \struck{[\ill{}} [l'une spécialisant l'autre],

\page{287}
où $A'^{\Gamma}$ n'est pas plat sur $S$, enfin où il n'y a pas de « bon » sous-schéma abélien des $A'^{\Gamma}$. Comme conséquence, on aura des exemples où 1°) la \uncertain{symétrie} est \uncertain{fausse} pour $X$ dans l'une ou l'autre \uncertain{direction} 2°) le \uncertain{thm} de \uncertain{continuité} est \uncertain{faux} pour $H^0(X, \Omega^1_X)$ et pour $H^1(X, \mathcal{O}_X)$ 3°) $\underline{\mathrm{Pic}}^{\tau}_{X/S}$ n'est pas plat sur $S$ 4°) Il n'y a pas de « bon » sous-schéma abélien de $\underline{\mathrm{Pic}}^{\tau}_{X/S}$, en particulier « la formation du $\mathrm{Pic}$ réduit n'est pas compatible avec la spécialisation ».

\subsection{Un homomorphisme divisible par $p$ sur la fibre spéciale seulement}

\page{288}
Soit $Y = \operatorname{Spec}(V)$, $V$ anneau de valuation discrète, $X$ schéma abélien sur $Y$, $y_0$, $y_1$ pt fermé resp.\ générique de $Y$, $X_0$, $X_1$ les deux fibres de $X$ …, $F$, $G$ deux sous-groupes de $X$, \uncertain{finis} sur $Y$,
\begin{itemize}
\item[(i)] $F_1 \not\subset G_1$, \struck{\ill{}}
\item[(ii)] $F_0 \subset G_0 \neq 0$
\item[(iii)] $pF_1 = G_1$ \quad ($p = \operatorname{car} k(V)$)
\end{itemize}
\note{les étiquettes sont écrites a), ii), iii) ; « iii) » surcharge une étiquette biffée, et l'indice de $F$ dans (iii) surcharge un autre chiffre}
\[
A = X/F , \qquad B = X/G , \qquad X \xrightarrow{\ \alpha\ } A , \quad X \xrightarrow{\ \beta\ } B \quad \text{hom.\ can.}
\]
Considérons $p\beta : X \to X/G$ (canoniques) alors $p\beta$ s'annule \struck{sur $pX$, donc} sur $F_1$ \add{en vertu de (iii)}, donc se factorise par $u : X/F \to X/G$

\begin{tikzcd}
& X \arrow[dl, "\alpha"'] \arrow[dr, "p\beta"] & \\
A \arrow[rr, "u"] & & B
\end{tikzcd}

\struck{Cependant} $u_1 \notin p\,\mathrm{Hom}(A_1, B_1)$, car \struck{\ill{}} \uncertain{si} \uncertain{plus} dans
\[
\mathrm{Hom}(X_1, X_1) \otimes \mathbb{Q} \simeq \mathrm{Hom}(A_1, B_1) \otimes \mathbb{Q} ,
\]
si \uncertain{on} \uncertain{avait} $u = pv$, on aurait $v = \beta$, or $\beta$ \uncertain{ne} \uncertain{provient} \ill{} d'un élément de $\mathrm{Hom}(A_1, B_1)$, puisque $F_1 \not\subset G_1$ \struck{[\ill{}, \ill{} (i)], \ill{} $F_1 = 0$ \ill{} \ill{}, \ill{} que contredit (ii)]}. Cependant $u_0 \in p\,\mathrm{Hom}(A_0, B_0)$, puisque $F_0 \subset G_0$.
\note{il écrit $p\,\mathrm{Hom}(A_0, F_0)$ ; le $F_0$ est sans doute un lapsus pour $B_0$}

Les conditions dites (i)(ii)(iii) peuvent être réalisées, en \uncertain{inégales} car., avec pour $X$ un \uncertain{produit} de deux \emph{courbes elliptiques}, en prenant $X_0$ sans pt d'ordre $p$, $X_1$ avec pt d'ordre $p$, et $F_1$ et $G_1$ deux s.-groupes de $X_1$, \struck{\ill{} \ill{}} \add{\uncertain{d'ordre} $p$, \ill{} $p$ \ill{} $k(y_1)$ et} distincts. Comme il n'y a

\page{289}
qu'un seul sous-groupe \struck{fini} de ${}_pX_0$ d'ordre $p$ sur $k(y_0)$, on aura $F_0 = G_0$, d'où \ill{} pour $F_1 \subset G_1$ et $pF = 0 \subset G$ \ill{} \ill{}, \uncertain{considérons} …,

$A$ et $B$ étant comme \uncertain{dessus}, \uncertain{considérons}
\[
C = A \times B' \qquad B' \text{ dual de } B
\]
et
\[
\underline{\mathrm{NS}}_{C/Y} \simeq \underline{\mathrm{NS}}_{A/Y} \times \underline{\mathrm{NS}}_{B/Y} \times \underline{\mathrm{Hom}}_{Y\text{-gr}}(A, B) .
\]
\note{après $\underline{\mathrm{NS}}$ un symbole biffé, puis une tache d'encre avant $\simeq$}
Alors $u$ correspond à une section $\bar u$ de $\underline{\mathrm{NS}}_{C/Y}$ [dont les premières deux composantes sont nulles], et sur la fibre \struck{\uncertain{générique}} spéciale, on a $\bar u_0 = p\bar v_0$, mais on ne peut avoir $\bar u = p\bar v$ car alors $\bar v$ \uncertain{nécessairement} aurait ses deux premières composantes nulles, donc proviendrait d'une section $v$ de $\underline{\mathrm{Hom}}_{Y\text{-gr}}(A, B)$, \ill{} on \uncertain{a} \uncertain{vu} n'était pas possible ….

\section{Structure de certains groupes divisibles}

\page{290}
\struck{\ill{}} Les \struck{\ill{}} groupes de torsion divisibles $M$ sont ceux qui sont isomorphes à des sommes directes de groupes $\mathbb{Q}_p/\mathbb{Z}_p$ ($p$ premier), ceux qui sont $p$-primaires sont ceux isomorphes à des \struck{\ill{}} \struck{\ill{} \ill{}} groupes $(\mathbb{Q}_p/\mathbb{Z}_p)^{(I)}$, où le cardinal de $I$ est bien déterminé ; c'est aussi la dimension de l'espace vectoriel sur $\mathbb{F}_p$
\[
{}_pM = \mathrm{Hom}(\mathbb{Z}/p\mathbb{Z}, M) .
\]
\note{la formule est encadrée ; à sa droite, deux lignes biffées : « \struck{Les \ill{} $M$ \ill{} groupes \ill{} \ill{}} »}
\struck{torsion $p$ primaires \ill{} \ill{} \ill{} \ill{} isomorphes}
Conditions équivalentes : (i) $M$ groupe de torsion $p$-primaire, et (ii) $M$ \uncertain{isomorphe} à un sous-groupe d'un \ill{} \ill{} \ill{} \ill{} ${}_pM$ est fini
\marginal{écrit en oblique dans la marge de gauche : modules \ill{} \ill{} \ill{} de \uncertain{Dieudonné}}
… $(\mathbb{Q}_p/\mathbb{Z}_p)^{\uncertain{n}}$, \struck{\ill{}} \ill{} \ill{} \ill{} \ill{} entier $\geq 0$. (iii) $M$ isomorphe à un groupe $F + (\mathbb{Q}_p/\mathbb{Z}_p)^{\rho}$, où $F$ est un $p$-groupe de torsion fini, \uncertain{et} $\rho$ un entier fini. Sous ces conditions, $\rho$ est défini de façon unique [c'est aussi le $\operatorname{rang}_{\mathbb{Z}/p\mathbb{Z}}$ \struck{$p$} ${}_pM^0$, où $M^0$ est le plus grand sous-groupe divisible de $M$] ainsi que la structure de $F$ [c'est $M/M^0$].

On peut le \uncertain{donner} \uncertain{sous} des façons différentes \uncertain{qui} \uncertain{correspondent} \ill{} : $M$ \ill{} \ill{} \ill{} de type fini sur $\mathbb{Z}_p$. On peut poser
\[
T_p(M) = \mathrm{Hom}(M, \mathbb{Q}_p/\mathbb{Z}_p) = \varprojlim_n \mathrm{Hom}({}_{p^n}M, \mathbb{Q}_p/\mathbb{Z}_p)
\]
\marginal{$T_p$ \emph{foncteur} \emph{contravariant} en $M$}
[$M \mapsto T_p(M)$ établit une \add{anti}équivalence de la catégorie des $M$ satisfaisant (i), et \uncertain{des} \uncertain{modules} de type fini sur $\mathbb{Z}_p$.]
On peut prendre aussi
\[
T^p(M) = \varprojlim_n {}_{p^n}M
\]
\marginal{$T^p$ \emph{foncteur} \emph{covariant} en $M$}
où pour $m \geq n$, on \uncertain{envoie} ${}_{p^m}M$ dans ${}_{p^n}M$ par $p^{m-n}$.

\page{291}
On a un isom.\ canonique
\[
T^p(M) = \mathrm{Hom}(T_p(M), \mathbb{Z}_p)
\]
On notera que a) $T^p$ n'est pas un foncteur exact, mais seulement \uncertain{exact} « \uncertain{modulo} groupes finis ». b) $T^p$ \uncertain{transforme} \ill{} les groupes finis \ill{} \uncertain{en} \uncertain{zéro} \ill{}
c) De façon précise, $T^p$ établit une \struck{\ill{}} équivalence des catégories « \emph{mod} les groupes finis », i.e.\ \uncertain{essentiellement} de \uncertain{celle} des \ill{} \ill{} \ill{} \ill{} de dim.\ finie sur $\mathbb{Q}_p$.

Soit $M \to N$ un \uncertain{homomorphisme} de groupes de la forme $(\mathbb{Q}_p/\mathbb{Z}_p)^n$, et $H$ un $p$-groupe fini. \emph{Conditions équivalentes} :
\begin{itemize}
\item[(i)] $T_p(N) \to T_p(M)$ injectif : \uncertain{conoyau} \uncertain{libre} [i.e.\ \uncertain{isom.} sur un facteur \uncertain{direct}]
\item[(i bis)] $T_p(N) \otimes H \to T_p(M) \otimes H$ injectif
\item[(i ter)] \struck{$\mathrm{Tor}_1(…, H)$}, $T_p(N)/pT_p(N) \to T_p(M)/pT_p(M)$ injectif
\item[(ii)] $T^p(M) \to T^p(N)$ surjectif
\item[(iii)] $M \to N$ \uncertain{induit} sur un quotient facteur direct
\item[(iv)] ${}_pM \to {}_pN$ est \emph{surjectif}
\end{itemize}

\section{Notes sur des feuillets dactylographiés}
\note{les pages 293, 296, 297 et 299 sont des notes de sa main écrites au verso ou en travers de feuilles dactylographiées sans rapport ; seules ses notes sont transcrites}

\page{293}
\marginal{en haut à droite, dans un cadre : $X$ famille du premier ordre \ill{} \uncertain{automatiquement} \ill{}}
$\underline{\mathrm{Pic}}^{\sigma\sigma}_{X/S}$ existe
\note{l'exposant de $\underline{\mathrm{Pic}}$ est douteux, peut-être $\tau$ ou $0$}
\[
\Updownarrow
\]
\[
\check V \xrightarrow{\ \mu\ } H^1(X_0, \underline{\mathcal{T}}_{X_0}) \xrightarrow{\ u\ } H^1(X_0, \mathcal{O}_{X_0}) \otimes t_{B_0} \simeq H^1(X_0, \mathcal{O}_{X_0}) \otimes H^1(A_0, \mathcal{O}_{A_0})
\]
$\operatorname{Im} u\mu \subset \operatorname{Im} v$, où
\[
v : H^1(A_0, \underline{\mathcal{T}}_{A_0}) \simeq t_{A_0} \otimes H^1(A_0, \mathcal{O}_{A_0}) \to H^1(X_0, \mathcal{O}_{X_0}) \otimes H^1(A_0, \mathcal{O}_{A_0})
\]
\note{la flèche $v$ est tracée verticalement, de bas en haut ; sa lettre est peu sûre}

$\underline{\mathrm{Pic}}_{X/S}$ plat sur $S$ en $x_0$
\[
\Downarrow
\]
\[
H^1(X_0, \mathcal{O}^*_{X_0}) \xrightarrow{\ \partial\ } \mathrm{Hom}\bigl(H^1(X_0, \underline{\mathcal{T}}_{X_0}), H^2(X_0, \mathcal{O}_{X_0})\bigr) \xrightarrow[\ \rho\ ]{} \mathrm{Hom}\bigl(\check V, H^2(X_0, \mathcal{O}_{X_0})\bigr)
\]
\note{au-dessus du dernier terme, il écrit « $V \otimes H^2(X_0, \mathcal{O}_{X_0})$ » relié par $\simeq$ ; le début du dernier terme surcharge un mot biffé illisible}
avec une flèche oblique $d : H^1(X_0, \mathcal{O}^*_{X_0}) \to H^1(X_0, \underline{\Omega}^1_{X_0})$ et une flèche verticale $c$, \uncertain{cup} $\vee$, de $H^1(X_0, \underline{\Omega}^1_{X_0})$ vers le $\mathrm{Hom}$ du milieu :
\[
\rho\partial(x) = 0 \iff \rho c(x) = 0
\]

N.B. Cette condition est automatiquement vérifiée si les $\xi^{(x_0)}$ \ill{} \ill{} $\underline{\mathrm{Pic}}^{\sigma}_{X_0/k}$ \struck{\ill{}}, mais il est douteux que $\underline{\mathrm{Pic}}^0_{X/K}$ soit pour autant plat sur $S$.
\marginal{écrit en oblique dans la marge de gauche : $\underline{\mathrm{Pic}}^0_{X/K}(k)$ \ill{} « \uncertain{divisible} »}

\page{296}
\[
\underline{\mathrm{Pic}}^{\tau}_{Y} \simeq \hat A_1 \times \hat B^{\Gamma}
\]
\[
\begin{aligned}
H^1(Y, \mathcal{O}_Y) &\simeq H^1(A_1, \mathcal{O}_{A_1}) \times H^1(B, \mathcal{O}_B) \simeq t_{\hat A_1} \times t_{\hat B_\Gamma} \\
H^1(X, \mathcal{O}_X) &\simeq H^1(A, \mathcal{O}_A) \times H^1(B, \mathcal{O}_B)
\end{aligned}
\]
\note{une flèche marquée $k$ descend sur $\mathcal{O}_{A_1}$ ; entre les deux lignes, une flèche verticale suivie de deux mots biffés illisibles ; une longue flèche courbe part de la première ligne vers la droite, et un trait la relie au terme encerclé « $H^1(B, \mathcal{O}_B) \otimes t_A$ », sous lequel $\simeq t_{\hat B}$}
\[
H^1(Y, \underline{\mathcal{T}}_Y) = H^1(Y, \mathcal{O}_Y) \otimes (t_A \times t_B) \longrightarrow H^1(Y, \mathcal{O}_Y) \otimes t_{A_1} \longleftarrow t_{\hat A_1} \otimes t_{A_1}
\]
\note{sous $H^1(Y, \mathcal{O}_Y)$, « $t_{\hat A_1} \times t_{\hat B}$ » ; au-dessus de $t_{A_1}$, « $t_A$ » relié par $\Vert$ ; sur la première flèche, deux mots : « \uncertain{naturellement} \uncertain{fonctorielle} \ill{} : $t_B$ »}
\[
\downarrow
\]
\[
H^1(X, \underline{\mathcal{T}}_X) = H^1(X, \mathcal{O}_X) \otimes (t_A \times t_B)
\]
\[
(t_{\hat A} \times t_{\hat B}) \otimes (t_A \times t_B) ,
\qquad t_{\hat B} \otimes t_A , \quad t_{\hat B} \otimes t_B
\]
\[
\hat A \leftarrow \hat A_1 , \qquad A \to A_1 , \qquad \hat A \leftarrow \hat A_1
\]
\[
t_A \otimes t_{\hat A} \rightleftarrows t_{A_1} \otimes t_{\hat A_1}
\]
\[
0 \to \mu_2 \to \underline{\mathrm{Pic}}^{*}_{Y} \to \underline{\mathrm{Pic}}^{\Gamma}_{X} \to \mu_2
\]
\note{sous $\underline{\mathrm{Pic}}^{*}_{Y}$ : « $\hat A_1 \times \hat B_{(2)}$ » ; sous $\underline{\mathrm{Pic}}^{\Gamma}_{X}$, relié par $\simeq$ : « $\hat A_1 \times \hat B_{(2)}$ », suivi d'un mot biffé ; l'indice de $\hat A$ dans ce dernier terme est surchargé}

\section{Contre-exemples pour recollement}
\note{titre de sa main, en haut à droite de la page 297}

\page{297}

\begin{tikzcd}[column sep=small, row sep=small]
& X_1 \amalg_Y X_2 & \\
X_1 \arrow[ur] & & X_2 \arrow[ul] \\
& Y \arrow[ul, "\varphi_1"] \arrow[ur, "\varphi_2"'] &
\end{tikzcd}

\note{sous le diagramme, « $\varphi_i$ immersions » ; à droite, un second diagramme vertical : $Y \rightrightarrows X_1 \amalg X_2 \to X_1 \amalg_Y X_2$, marqué d'un astérisque}

\emph{exemple} : $X_1 = X_2 = \operatorname{Spec} k[x, y]$ -- origine
\[
\text{\struck{$Y = \operatorname{Spec} k[t]$ -- origine}} = \operatorname{Spec} k[t, t^{-1}]
\]
\note{sous la ligne biffée, la lettre $Z$}
$\varphi_1$ \uncertain{défini} par \struck{\ill{}} $Y \to X$ défini \uncertain{par} \uncertain{les} équation\supplied{s} : $x = \struck{\ill{}}\, t$, $y = 0$
\struck{$\varphi_2$ \ill{} $Z \to X$ \ill{} $x = 1/t$ \ill{}}
\struck{$\varphi_2\psi = \varphi_1\psi$ … $\psi$}
$x = 1/t$, $y = 0$
\note{un passage de quatre lignes est biffé de longs traits ; dans la marge, une petite croix d'axes}
$X' = X_1 \amalg X_2 \to$ \ill{}

$X = X_1 \amalg_Y X_2 \to$ \ill{}

$\Gamma(X, \mathcal{O}_X) =$ \struck{\ill{}} …, $\Gamma(X')$ formé des $(f_1, f_2)$ tels que $f_1$ \struck{\ill{}} $\circ\, \varphi_1 = f_2 \circ \varphi_2$.
Or $\Gamma(X_i, \mathcal{O}_{X_i}) = k[x, y]$, et \ill{} \ill{} $f_1(t, 0) = f_2(1/t, 0)$

\page{299}
\struck{\ill{} \ill{} \uncertain{exige}} \add{\uncertain{dire}} $g_1(t) = g_2(1/t)$ \uncertain{en} \struck{\ill{} \ill{}}
\uncertain{g\supplied{al}e} \uncertain{une} \uncertain{constante}. Donc les

\uncertain{fonctions} de $Y$ \uncertain{ne} \uncertain{sont} pas séparés \ill{}

On \ill{} \uncertain{aussi} \ill{} \ill{} \uncertain{variétés} par le groupe $\mathbb{Z}/2$
[\uncertain{opérant} \uncertain{sur} $X_1 \amalg X_2$ par \uncertain{symétrie}, \marginal{écrit en oblique dans la marge : \uncertain{connexité} ; \uncertain{irréductible} …}
et sur $Y$ par $t \mapsto 1/t$] \ill{}

\ill{}
\[
(X_1 \amalg_Y X_2)/\mathbb{Z}_2
\]
\ill{} $X_1$ \ill{} \uncertain{affine}.
\[
\uparrow
\]
\[
X_1 \amalg X_2/\mathbb{Z}/2 \simeq X_1
\]
\ill{} \ill{} $X_1 \amalg_Y X_2/\mathbb{Z}_2$
puis que $X_1 \amalg_Y X_2$ le \uncertain{serait} \ill{}.

Il \uncertain{faut} \uncertain{cependant} vérifier que $H$ \struck{\ill{}} \add{\ill{}} \struck{\ill{}} de $X_1 \amalg_Y X_2$ \uncertain{par} $\mathbb{Z}/2$ \ill{} \ill{} \ill{} \ill{} affines c'est facile. \add{\ill{} \ill{}} \ill{} \ill{} \ill{} \ill{} $D$ \ill{} $\mathfrak{X}$, \ill{} \ill{} \ill{} \ill{} \ill{}, \uncertain{recollés} $(\mathfrak{X} - D) \amalg (\mathfrak{X} - D)$ …
\note{l'argument se poursuit au-delà de la page 299 ; la page 300 est sans rapport et la suite, s'il y en a une, est hors de ce lot}

\end{document}
